\section{Profligacy of thought}

Let's say a person has a thought every two seconds -- that seems about right, although I am sure that thoughts just race through a large number of people, and they are slow to form in a much smaller number. Assuming 16 waking hours per day, that comes to about 30,000 thoughts per day, or 11,000,000 thoughts per year. So in 65 years of thinking, there are 715,000,000 thoughts, or 715 mega-thoughts in a lifetime.

Another way to consider it, is that with the current world population at 6,800 million, that means the entire thought output of the human race is around is 74,800,000,000,000,000 thoughts, or 74.8 peta-thoughts (74.8 quadrillion short form or 74.8 billiard long form), per year.

Someone once suggested that if a million monkeys typed letters randomly, they eventually would create the entire works of Shakespeare. Well, let's investigate that. The 74.8 peta-thoughts in the year 2009 certainly did not create the works of Shakespeare; perhaps some good books and plays, but nothing quite that memorable.

The number of humans who have every lived\footnote{\url{http://en.wikipedia.org/wiki/World\_population}} is perhaps 100 billion (100,000 million). So the entire thought output since the beginning of time is about 7,480,000,000,000,000,000,000,000,000 thoughts, or 7,480 yotta-thoughts.

Is it unreasonable to suppose that with that many thoughts, even randomly, every human problem should have been solved, and every scientific theory understood? Why are we not even close? Let's examine what we think about in Table \ref{tab:levesnegativethoughts}:

\begin{table}[h]\small\centering
\begin{tabular}{rl}\toprule
Negative thoughts & Worry, anxiety, fear, concern about what others a thinking \\

Fantasies & Fantasies about sex, glory, revenge, conquest, future plans \\

Daily life & What should I eat, what should I do today, where should I go\\\bottomrule
\end{tabular}
\caption{Levels of negative thoughts}
\label{tab:levesnegativethoughts}
\end{table}

These take up most of the day. And we see that thoughts are not random, since they come in certain sequences. The process is mechanical, so a given thought more or less determines the next thought, and the next in a constricting chain, until boredom or an external stimuli starts a new sequence.

These are at the level of doxa. Even at the level of dianoia, thought is still mechanical. People prefer to repeat theories, slogans, phrases that have already been said, over and over again. Or conversation turns argumentative rather than constructive, since the participants prefer to stand their ground rather than be persuaded by facts or logical arguments.

Again, none of this is random, so we can't count on randomness to generate much in the way of novelty. Then, there is the price of admission to the highest levels of science, theology, history, as well as the creative arts. It is limited to the very few of superior intelligence with access to the proper education. Even given that, there is very little novelty. The number of thinkers who have made major contributions to thought is quite minuscule in relation to the total world population, past and present.

Beyond that is episteme, or true thought; that is, thought that concentrates on the world of ideas, beyond sense perception, beyond the personal concerns of the moment. This requires discipline and spiritual training. Then, among those with such insights, there are even fewer who can communicate those insights and command allegiance to them. The number of great founders of religious traditions, and the next tier of mystics and saints below them, is indeed quite small.


\flrightit{Posted on 2010-05-05 by Cologero}

\begin{center}* * *\end{center}

\begin{footnotesize}
\begin{sffamily}
\texttt{Katerina on 2010-05-05 at 23:10 said:}

There are original thinkers (ie: The Greeks)on philosophy, religion, the origins of the universe, etc, etc, and then there are ``the others''. You are in the first category Cologero. I think that with so many human thoughts, even randomly, every problem has been solved, although sometimes the person has not understood what they have in their head, and does not put these thoughts into words. They may think others will assume they are crazy, or worse. Then there are those who don't care, and expound upon myriad theories until one sticks. I believe we are all born knowing the answer to everything in our subconscious, the collective knowledge of millenia is imprinted in our DNA, but intellectualizing these thoughts is the problem for many.


\hfill

\end{sffamily}
\end{footnotesize}

